%-------------------------
% Based off of: https://github.com/sb2nov/resume
% License : MIT
%------------------------

\documentclass[letterpaper,12pt]{article}

\usepackage{latexsym}
\usepackage[empty]{fullpage}
\usepackage{titlesec}
\usepackage{marvosym}
\usepackage[usenames,dvipsnames]{color}
\usepackage{verbatim}
\usepackage{enumitem}
\usepackage[hidelinks]{hyperref}
\usepackage{fancyhdr}
\usepackage[russian, english]{babel}
\usepackage{tabularx}
\input{glyphtounicode}


%----------FONT OPTIONS----------
% sans-serif
% \usepackage[sfdefault]{FiraSans}
% \usepackage[sfdefault]{roboto}
% \usepackage[sfdefault]{noto-sans}
% \usepackage[default]{sourcesanspro}

% serif
% \usepackage{CormorantGaramond}
% \usepackage{charter}


\pagestyle{fancy}
\fancyhf{} % clear all header and footer fields
\fancyfoot{}
\renewcommand{\headrulewidth}{0pt}
\renewcommand{\footrulewidth}{0pt}

% Adjust margins
\addtolength{\oddsidemargin}{-0.5in}
\addtolength{\evensidemargin}{-0.5in}
\addtolength{\textwidth}{1in}
\addtolength{\topmargin}{-.5in}
\addtolength{\textheight}{1.0in}

\urlstyle{same}

\raggedbottom
\raggedright
\setlength{\tabcolsep}{0in}

% Sections formatting
\titleformat{\section}{
	\vspace{-4pt}\scshape\raggedright\large
}{}{0em}{}[\color{black}\titlerule \vspace{-5pt}]

% Ensure that generate pdf is machine readable/ATS parsable
\pdfgentounicode=1

%-------------------------
% Custom commands
\newcommand{\resumeItem}[2]{
	\item\small{
		{\ifx&#1&\else\textbf{#1}: \fi#2 \vspace{-2pt}}
	}
}

\usepackage[T2A]{fontenc}
\usepackage[utf8]{inputenc}
\usepackage{cmap}

\newcommand{\resumeSubheading}[4]{
	\vspace{-1pt}\item
	\begin{tabular*}{0.97\textwidth}{l@{\extracolsep{\fill}}r}
		\textbf{#1} & #2 \\
		\ifx&#3&\else\textit{\small#3} & \textit{\small #4} \\ \fi
	\end{tabular*}\vspace{-5pt}
}


\newcommand{\resumeSubSubheading}[2]{
	\item
	\begin{tabular*}{0.97\textwidth}{l@{\extracolsep{\fill}}r}
		\textit{\small#1} & \textit{\small #2} \\
	\end{tabular*}\vspace{-7pt}
}

\newcommand{\resumeProjectHeading}[2]{
	\item
	\begin{tabular*}{0.97\textwidth}{l@{\extracolsep{\fill}}r}
		\small#1 & #2 \\
	\end{tabular*}\vspace{-7pt}
}

\newcommand{\resumeSubItem}[2]{\resumeItem{#1}{#2}\vspace{-4pt}}

\renewcommand\labelitemii{$\vcenter{\hbox{\tiny$\bullet$}}$}

\newcommand{\resumeSubHeadingListStart}{\begin{itemize}[leftmargin=0.15in, label={}]}
	\newcommand{\resumeSubHeadingListEnd}{\end{itemize}}
\newcommand{\resumeItemListStart}{\begin{itemize}}
	\newcommand{\resumeItemListEnd}{\end{itemize}\vspace{-5pt}}
\newcommand{\ulhref}[2]{\href{#1}{\underline{#2}}}
%-------------------------------------------
%%%%%%  RESUME STARTS HERE  %%%%%%%%%%%%%%%%%%%%%%%%%%%%


\begin{document}

	\begin{center}
		\textbf{\Large \scshape Усатов Павел} \textbf{\Large $|$ DevOps Engineer} \\ \vspace{1pt}
		\small +7 (912) 712-83-52 $|$
		\href{https://t.me/Usatov_Pavel}{\underline{Telegram}} $|$
		\href{mailto:pvusatov@edu.hse.ru}{\underline{pvusatov@edu.hse.ru}} $|$
		\href{https://github.com/UsatovPavel}{\underline{Github}} $|$
		\href{https://usatovpavel.ru/default}{\underline{\textbf{Сайт}}}
	\end{center}

	%-----------EDUCATION-----------
	\section{Образование}
	\resumeSubHeadingListStart
	\resumeSubheading
	{ВШЭ СПБ}{Санкт-Петербург}
	{Прикладная математика и информатика}{2023-2027}
    \item \ulhref{https://usatovpavel.ru/jvm\#activities}{Интенсивы ШАД Яндекса}: Agents Week и AI Agents Security Week 
	\resumeSubHeadingListEnd

	%-----------EXPERIENCE-----------
	\section{Опыт}
	\resumeSubHeadingListStart
	\resumeSubheading
	{\ulhref{https://github.com/UsatovPavel/riid}{Утилита RIID для внутреннего облака VK}}{январь 2026 - июнь 2026}
	{Daemon доставки образов по p2p, независимый от container engine}{Курсовая, научрук - SRE из VK}
	\resumeItemListStart
	\resumeItem{}{Развернул в k8s: DaemonSet, ConfigMap, Secret, Service, local-path storage}
	\resumeItem{}{CI в GitHub Actions: minikube и Dragonfly в раннере, релиз по тегу в Packages}
	\resumeItem{}{Helm-чарт riid-observer: VictoriaMetrics, Grafana, kube-state-metrics}
	\resumeItem{}{Тестирование: 12 нод, 100 образов от 1 МБ до 5 ГБ: на 20\% быстрее Podman}
	\resumeItemListEnd
	\textbf{Технологии}: Kubernetes, Helm, Docker, Podman, Prometheus, Grafana, Yandex Cloud

	\resumeSubheading
	{Стажировка Backend Т-Банк}{апрель 2026 - июль 2026}
	{Продуктовая команда документооборота}{Kanban}
	\resumeItemListStart
	\resumeItem{}{Версионировал эндпоинты и катил через GitLab CI, не ломая мобильные клиенты}
	\resumeItem{}{Алерты на БД, оптимизация запроса, установка расширений PostgreSQL}
	\resumeItem{}{Написал модуль команды в общем репозитории интеграции b2b, что позволило перейти другим командам. Проверил e2e на QA-контуре}
	\resumeItem{}{Разбирал баги генерации документов по логам в ELK}
	\resumeItemListEnd
	\textbf{Технологии}: GitLab CI, PostgreSQL, Kafka, Docker, ELK, Scala

	\resumeSubHeadingListEnd

	%-----------PROJECTS-----------
	\section{Проекты}
	\resumeSubHeadingListStart
	\resumeSubheading
	{\ulhref{https://github.com/UsatovAI/ansible-deploy}{ansible-deploy} $|$ \normalfont{\emph{\ulhref{https://github.com/UsatovAI/claude-web-console}{claude-web-console}}}}{2026}
	{Автодеплой веб-консоли к headless Claude Code на VPS}{Индивидуальный проект}
	\resumeItemListStart
	\resumeItem{}{Ansible-роли с тегами и секретами в ansible-vault; прод и тест - разные группы}
	\resumeItem{}{systemd без root с ProtectSystem=strict, certbot с renewal-хуком, SNI на nginx}
	\resumeItemListEnd
	\textbf{Технологии}: Ansible, ansible-vault, systemd, Nginx, certbot/Let's Encrypt, Bash

	\resumeSubheading
	{\ulhref{https://github.com/UsatovPavel/PRAssign}{PRAssign} $|$ \normalfont{\emph{сервис управления code review}}}{ноябрь 2025 - январь 2026}
	{Go REST API и асинхронный пайплайн на Kafka}{Индивидуальный проект}
	\resumeItemListStart
	\resumeItem{}{Compose для сборки/тестов/нагрузки, Nginx балансирует реплики с повтором по 5xx}
	\resumeItemListEnd
	\textbf{Технологии}: Go, Nginx, Kafka, PostgreSQL, Docker Compose, k6, golangci-lint

	\resumeSubheading
	{\ulhref{https://github.com/hse-project-Java-2025}{TimeTamer}}{2025}
	{Java-сервер и Android-приложение календаря}{Командный проект}
	\resumeItemListStart
	\resumeItem{}{Докеризовал сервер: Dockerfile и compose с PostgreSQL, кэш слоёв сборки}
	\resumeItemListEnd
	\textbf{Технологии}: Docker, Docker Compose, PostgreSQL, Java, Spring Boot
	\resumeSubHeadingListEnd

	%-----------SKILLS-----------
	\section{Навыки}
	\textbf{Kubernetes и контейнеры}{: DaemonSet, ConfigMap, Helm, Docker, Compose, Podman} \\
	\textbf{CI/CD и конфигурация}{: Git, GitLab CI, GitHub Actions, Ansible, Make, Bash, Gradle} \\
	\textbf{Мониторинг и ОС}{: Prometheus, VictoriaMetrics, Grafana, ELK, Linux, systemd, Nginx, TLS} \\
	\textbf{Языки и данные}{: Go, Python, Java, Scala, PostgreSQL, Kafka, gRPC}

%-------------------------------------------
\end{document}
